% Goldman 生益 TP reconciled: 217.6 headline (later Sina repost, 45x 2027E PE) per footnote below;
% 127.4/146.3 (older note, 35x 2027E PE) superseded. 217.6 (Goldman/Sina, 45x 2027E PE) is the sole
% public-sentiment ceiling; 195 (Citi) sits at ~82% of the 150-205 band and is NOT a ceiling.
% Applied consistently across ch06/ch07/ch08/ch09 (ch09 line 670 uses only 217.6 as the ceiling).
\section{公开研究情绪、共识聚类与多空平衡}

本章讨论市场共识，而不是复述卖方观点。公开研究情绪可以帮助识别哪些逻辑已经被充分交易、哪些变量仍有分歧，但最终投资动作仍要回到AStock自己的证据、估值和风险框架。

\begin{sourcequalitybox}[为何本章不是完整街头一致预期]
本次语料已取得多份原始券商PDF和官方IR记录，但没有覆盖Wind/Choice级别的全部机构预测数据库、全部全球投行原文和连续目标价历史。因此本章仍称为``公开可得研究情绪''，而不是完整街头一致预期。这些观点可作为证据，但不是报告自身声音。
\end{sourcequalitybox}

公开研究情绪的价值在于识别市场共识和盲点，而不是替代本报告判断。当前卖方材料基本一致看多AI PCB、CCL升级和封装基板机会，但不同机构使用的证据强度不一致：有些来自原始PDF和模型表，有些来自转载摘要或公开预览。因此，本章只把卖方观点作为“市场正在相信什么”的输入，再由AStock结合官方交付、估值和客户链重新定价。

\section{共识聚类图}

\begin{exhibitbox}[图表7-1：公开卖方材料的共识聚类（What public sell-side materials agree on）]
\small
\begin{tabularx}{\textwidth}{L{3cm}C{1.8cm}X X}
\toprule
\textbf{Cluster} & \textbf{Tone} & \textbf{Source support} & \textbf{AStock treatment} \\
\midrule
PCB半导体化 & \cellcolor{green!18}Bullish & 国金、国盛、高盛、摩根大通公开材料；JPM胜宏Reportify/新浪可见正文补充。 & 核心论点，但绑定估值纪律；JPM目标和预测只作H股/公开情绪输入。 \\
CCL价值提升 & \cellcolor{green!18}Bullish & 国盛、高盛、花旗、申万宏源公开材料；高盛/花旗生益新浪可见正文补充。 & 材料链主线，跟踪M8/M9价格和交期；外资目标价不替代AStock合理区间。 \\
封装基板期权 & \cellcolor{yellow!18}Constructive & 国海、西部、申万宏源公开材料。 & 结构性期权，关注认证和良率。 \\
设备耗材 & \cellcolor{yellow!18}Constructive & 东吴机械和设备总结材料。 & 二阶CAPEX弹性，不作为最稳核心资产。 \\
国产算力链 & \cellcolor{yellow!18}Constructive & 国盛、申万宏源公开材料。 & 长期主题，短期证据不如北美AI链直接。 \\
\bottomrule
\end{tabularx}
\sourcenote{图表7-1：高盛/花旗/摩根大通/国金/国盛/申万宏源/国海/西部/东吴公开材料及新浪可见转载正文；AStock内部估值与客户链证据。聚类的AStock覆盖标的映射：PCB半导体化→沪电002463/胜宏300476；CCL价值提升→生益600183/华正603186；封装基板→深南002916；设备耗材→芯碁688630；国产算力链→参照扩大覆盖标的。}
\end{exhibitbox}

共识最强的地方，也是估值最容易透支的地方。PCB半导体化和CCL价值提升已经成为主流叙事，意味着这些方向需要更强的业绩验证才能继续上修；封装基板、设备耗材和国产算力链仍有结构性机会，但短期证据不如北美AI链和高端PCB直接。

两个``Bullish''聚类的因果逻辑并不对称。PCB半导体化共识最强，证据来自具名券商PDF/转载正文与官方交付，但正因如此也最容易被price in。对照具体数字：沪电股份（002463）AStock合理价值区间为130--160元（见估值章图表8-1），高盛/新浪转载目标价142元落在区间内部（约第40分位），并未突破上沿，说明即便龙头的卖方目标价也已被AStock纪律区间吸纳，后续上行只能依赖订单和毛利率的同步兑现而非估值重估；胜宏科技（300476）AStock合理价值区间为290--415元，但卖方仅给出摩根大通H股600港元目标价，因A/H证券类别和预测口径不可比，不能直接与A股人民币区间换算为``差距''，其定价透支程度须回到隐含净利门槛（见估值章图表8-6）而非目标价差。CCL价值提升共识的证据强度则更依赖一个尚未闭合的变量——M8/M9等级CCL的ASP与交期，该变量正处于本报告``不可标记完成''的证据边界上（见风险章风险框架），因此CCL链的多头逻辑虽方向正确，触发点仍未到齐。封装基板、设备耗材和国产算力链属于``Constructive''而非``Bullish''：方向获得多数机构认可，但短期证据密度低于北美AI链和高端PCB，更适合作为结构性期权而非最稳核心资产。

值得强调的是，``共识看多''本身是多头假设的一个风险项，而非单纯利好。把上述两个Bullish聚类的卖方目标价与AStock自身估值输出对照——以生益科技为例，高盛转载目标价217.6元高于AStock合理价值区间150--205元上沿（见估值章，参考价180.15元），构成唯一真正突破纪律区间上沿的外部目标价，也是公开情绪的``天花板参照''；而花旗转载目标价195元落在区间内部（约第82分位），仅表明其仍处AStock可接受区间高位，并非突破上沿。两者性质不同：高盛217.6元属外部情绪给出、超出报告纪律区间的``天花板''，应被视为兑现风险监测器而非追加买入信号；花旗195元则说明即便较温和的外部预期也已逼近区间上沿。二级市场实际持仓与资金流的对应判断见二级市场章。

已归档的摩根大通/高盛/花旗转载正文显示，公开情绪层的预测输入更完整：摩根大通给出胜宏H股600港元目标价和2026--2028年收入/净利高增长情景，高盛和花旗分别给出生益科技217.6元和195元目标价及高端CCL景气假设（其中仅高盛217.6元突破AStock合理价值区间150--205元上沿，花旗195元仍处区间内部，见上文）\footnote{语料中存在两份高盛生益材料：较早的归档材料目标价127.4元（35倍2027E PE），较晚的新浪转载材料目标价由146.3元上调至217.6元（45倍2027E PE）。本章以较晚的217.6元作为公开情绪上限引用，较早的127.4元视为已被替代。}。但这些材料仍不是原始PDF，也不能把NVIDIA、Google、Rubin等平台叙事直接转成具名客户收入拆分。

对转载页内嵌图片进行了中文OCR\footnote{OCR方法学细节（中文模型、转载图像下载流程）见方法学附录。}，补出摩根大通胜宏收入假设表的机器可读线索：2026--2028E总收入约354/600/842亿元，MLPCB约157/282/407亿元，HDI约161/278/393亿元；花旗生益图片还补充了电子玻纤布1080/2116/7628系列年内涨幅约95\%/91\%/60\%的材料价格线索。OCR仍是转载图像的补充证据（OCR还原值，非原始PDF，待交叉核对），不等于原始PDF。

\section{多空平衡}

多空平衡部分把共识拆成可证伪命题。多头的核心是假设AI平台升级提高板级价值量，并且高端材料/产能持续紧缺；空头的核心不是否认AI需求，而是质疑资本开支节奏、材料涨价持续性、技术路线变化和2027--2028年产能过剩。后续投资跟踪应围绕这些命题，而不是围绕主题热度本身。

\begin{dashboardbox}[图表7-2：多空命题与盲点（What the market may be missing）]
\small
\begin{tabularx}{\textwidth}{X X}
\toprule
\textbf{Bull case evidence} & \textbf{Bear case / blind spot} \\
\midrule
AI服务器平台升级提高PCB层数、材料等级和背板/中板需求。 & AI资本开支可能出现消化期，订单可见度比叙事更重要。 \\
M9 CCL和高端材料成为信号完整性瓶颈。 & 材料涨价若反转，PCB厂和材料厂利润传导会重新分配。 \\
先进封装模糊PCB、载板和封装边界。 & 玻璃基板、光互联或新架构可能改变PCB价值池。 \\
扩产和技术通胀带动设备耗材订单。 & 2027--2028年可能出现产能过剩和订单波动。 \\
\bottomrule
\end{tabularx}
\sourcenote{图表7-2：四组多空命题由AStock基于本章``共识聚类''与多空平衡框架综合整理，证据来源同图表7-1；命题绑定具名标的的验证锚见下文。}
\end{dashboardbox}

把上述四组多空命题落到具名标的与可监测触发器上，才能从``命题清单''变成``监测清单''。多头第1项（AI服务器平台升级）的验证锚是胜宏（300476）2026E MLPCB收入是否触及摩根大通OCR口径的157亿元以及沪电（002463）高速交换机订单环比，二者同步走强则维持进攻观察权重，任一走弱则把胜宏弹性权重下修；多头第2项（M9 CCL）对应生益（600183）与华正（603186）的M8/M9等级ASP与交期，是``不可标记完成''证据边界，触发前不追加材料链仓位。空头第3项（封装/基板/光互联新架构）主要威胁深南（002916）FC-BGA良率与扩产节奏，早期信号是认证进度而非主题热度；空头第4项（2027--2028产能过剩）的早期信号是设备订单同比与深南载板良率，一旦设备订单同比转负或良率爬坡滞后，应触发降仓位而非依赖主题热度。换言之，多空命题的解算顺序决定持仓动作：先看订单与毛利率（多头第1项），再看材料ASP与交期（多头第2项），最后才看产能过剩信号（空头第4项）。

如果中报显示订单和毛利率同步兑现，多头叙事会继续占优；如果只看到资本开支增加、收入增长但毛利率和现金流走弱，则空头叙事会获得更多证据。这个框架也决定了本报告的投资动作：先验证，再加权，而不是因为共识强就追价。

\section{扩充标的情绪画像}

前述共识聚类与多空平衡主要围绕核心6家（沪电/深南/胜宏/生益/华正/鹏鼎）展开。为使情绪覆盖度与全报告20只标的覆盖范围匹配，本节把情绪分析延伸至14只扩充标的，按投资分层组织为四类：Core Candidate（3家，可投资但需更宽估值带）、Watchlist（9家，仅信号跟踪不设目标价）、Demand Indicator（2家，需求侧情绪温度计）、HK Anchor（1家，港股CCL锚）。每类标的的情绪颗粒度与其证据密度和投资角色匹配——Core Candidate按段分析，Watchlist用综合情绪对比表浓缩呈现，Demand Indicator和HK Anchor则承担特定维度的参照功能。

\begin{sourcequalitybox}[扩充标的情绪数据边界]
14只扩充标的的情绪信号取自公开东方财富龙虎榜、融资融券、股东增减持、股权质押、限售解禁、资金流向、基金持有等代理指标，以及公司公告（风险提示、澄清公告、监管函）和投资者关系活动记录。这些信号的密度低于核心6家（例如：扩充标的无完整300日融资融券长窗、北向持股冻结于2024-08-16页面壳、机构持仓为公共持有人行而非重仓排名），因此情绪判断的置信度低于核心标的，仅作信号监测和交叉验证，不作为独立交易依据。
\end{sourcequalitybox}

\subsection*{Core Candidate情绪画像（3家）}

Core Candidate三家（东山精密、景旺电子、芯碁微装）均具备清晰的AI链卡位证据，且市场关注度和交易活跃度已进入主流资金视野。三家的共性是估值均含显著AI溢价，但内部人动作、监管互动和热度结构差异较大，决定了各自的情绪安全边际。

\paragraph{东山精密（002384）} 估值处于历史高位区间——按2026Q1年化EPS测算动态PE约112倍，按2025全年净利计TTM PE约227倍，显著高于核心6家均值，定价中已充分计入索尔思光电并表与1.6T光模块垂直一体化预期。内部人动作方面，公开增减持记录未见董监高大额减持，但2026年以来股东户数呈扩散趋势，筹码有从集中向分散迁移的迹象。监管互动层面，公司于2026年6月密集发布光芯片/光模块资本开支公告（索尔思USD 1.2bn扩产），信息披露节奏与股价上行高度同步，需关注后续是否触发交易所关注函。市场热度极高：近90日7次登上龙虎榜（扩充标的中仅次于金安国纪），融资余额约141.39亿元（扩充标的中最高）、占流通市值3.97\%，单日成交额超200亿元属常态。情绪温度计读数偏热，定价对索尔思整合节奏和Google ASIC客户链验证高度敏感，任何不及预期都可能触发剧烈回调。

\paragraph{景旺电子（603228）} 估值相对克制——PE-TTM约69倍、PB约6.0倍，低于沪电/胜宏/深南一线均值，呈现"AI溢价未充分但仍贵于纯制造可比"的中间状态。内部人动作方面，公开增减持数据未见董监高近期大额减持，但公司2026Q1出现"增收不增利"剪刀差（营收+16.41\%、归母-28.37\%、扣非-18\%、现金流-52\%），与当前估值水平形成基本面-情绪张力。监管互动平稳，未见风险提示或监管函，但泰国基地投产进度和1.6T光模块PCB出货量是市场持续追问的IR焦点。市场热度中等：近90日未登上龙虎榜（扩充标的中少数未上榜的公司之一），融资余额约26.25亿元、占流通市值3.32\%，基金持有约998行、合计持股市值约179.12亿元。情绪温度计读数温偏凉，属于"有故事但未炒透"的标的，关键观察点是Q2/Q3利润端能否修复正增长——修复则估值折价收窄，否则"增收不增利"叙事会压制情绪。

\paragraph{芯碁微装（688630）} 估值已显著透支——2026E PE约127倍、2027E约91倍，现价502元已超卖方目标价区间（243--326.82元）上限约54\%，是扩充标的中"现价vs卖方目标价倒挂"最严重的一只。内部人动作方面，公开增减持记录未见董监高近期减持，但2026年6月26日H股挂牌为重大情绪事件——H股定价若显著低于A股，将对A股高估值构成反向牵引。监管互动正常，投资者关系活动记录披露了WLP2000晶圆级封装设备量产测试进展，但客户名称和订单金额均未披露，属"needs-IR"证据边界。市场热度偏高：近90日2次上榜、全统计窗口5次登龙虎榜，融资余额约10.94亿元、占流通市值1.65\%，基金持有约992行、合计持股市值约88.88亿元。情绪温度计读数极热但边际有松动迹象——30日主力资金净流仅-0.16亿元（接近平盘），且估值已超越所有卖方目标价，属于"情绪驱动+故事期"定价，需严守触发器纪律，一旦H股定价偏低或OCF延续性证伪即触发估值压缩。

\subsection*{Watchlist情绪对比表（9家）}

Watchlist共9只标的，共性是PE极端（或失效）、当期AI收入占比极低、或公司主动澄清无AI业务。对这类标的进行逐段情绪分析效率低且容易放大噪音，因此采用综合情绪对比表浓缩呈现，列出估值分位、内部人动作、监管/澄清信号、龙虎榜频次、融资融券强度、资金流向等关键情绪维度，便于快速扫描和预警。

\begin{exhibitbox}[图表7-3：扩充Watchlist情绪信号综合对比表]
\scriptsize
\begin{tabularx}{\textwidth}{L{1.6cm}C{1.5cm}L{1.5cm}L{1.8cm}C{0.8cm}C{1.2cm}C{1.3cm}}
\toprule
\textbf{标的} & \textbf{估值水位} & \textbf{内部人动作} & \textbf{监管/澄清} & \textbf{LHB} & \textbf{融资余额/占比} & \textbf{30日主力净流} \\
\midrule
铜冠铜箔 \newline 301217 & \cellcolor{red!12}PE-TTM ~287x \newline 极端高位 & 未披露大额增减持 & 未见监管函；\newline HVLP5/载体铜箔未量产 & 6次 & 13.98亿 \newline (0.93\%) & -11.38亿 \\
东材科技 \newline 601208 & \cellcolor{red!12}PE ~200x \newline 历史高位 & \cellcolor{red!12}董事长减持0.29\% & 未见监管函；\newline Rubin预期为needs-IR & 1次 & 31.88亿 \newline (4.22\%) & -22.53亿 \\
宏昌电子 \newline 603002 & \cellcolor{red!20}PE >4,000x \newline 失效 & \cellcolor{red!12}财务负责人减持 & \cellcolor{red!12}公司紧急风险提示；\newline 澄清AI-PCB未转利润 & 3次 & 非两融 & -2.20亿 \\
宏和科技 \newline 603256 & \cellcolor{red!20}PE ~750x \newline 叙事溢价极值 & \cellcolor{red!12}董事长/董秘 \newline 5月高位减持 & \cellcolor{red!12}公司6/17提示\newline 扩产不确定 & 5次 & 非两融 & +1.62亿 \\
中材科技 \newline 002080 & \cellcolor{yellow!18}PE ~94x \newline 偏高 & 未披露大额减持 & \cellcolor{red!12}两次风险提示；\newline AI电子布仅2.08\% & 5次 & 20.96亿 \newline (1.52\%) & -10.85亿 \\
中国巨石 \newline 600176 & \cellcolor{green!12}PE ~42x \newline 周期高点 & 未披露大额减持 & \cellcolor{yellow!18}公司6/17澄清\newline AI敞口0\% & 1次 & 34.77亿 \newline (1.62\%) & -19.14亿 \\
金安国纪 \newline 002636 & \cellcolor{red!12}PE ~107--253x \newline 伪AI泡沫 & 未披露大额减持 & \cellcolor{red!20}两次澄清无AI敞口；\newline 监管异常波动关注 & 9次 & 18.77亿 \newline (2.71\%) & -9.05亿 \\
容大感光 \newline 300576 & \cellcolor{red!12}PE ~165x \newline 3年分位96.83\% & 未披露大额减持 & \cellcolor{red!12}深交所监管函；\newline 募投理财超期+信披违规 & 1次 & 6.31亿 \newline (4.87\%) & -3.10亿 \\
\rowcolor{blue!6}
建滔积层板 \newline 01888.HK & 股价年内+570\% \newline （港股锚） & 未披露大额减持 & 港股无A股式监管函；\newline 涨价>50\%已公告 & N/A & 港股不适用 & 南向8日\newline +96.53亿港元 \\
\bottomrule
\end{tabularx}
\vspace{0.3em}
\sourcenote{图表7-3：数据来源——估值水位取自本报告估值章图表8-25及各标的evidence card；内部人动作取自东方财富股东增减持数据库（RP\_SHARE\_HOLDER\_INCREASE）及公司公告；监管/澄清取自深交所/上交所公告、公司风险提示公告与监管函；LHB取自东方财富龙虎榜数据（近90日窗口：2026-03-22至2026-06-20）；融资融券取自Eastmoney RPTA\_WEB\_RZRQ\_GGMX（2026-06-17/18快照，非两融标的标记为"非两融"）；30日主力净流取自东方财富日频资金流向代理指标。建滔积层板为港股，A股端点数据不适用，替换为港股通南向净买入数据。颜色标记：绿=估值或情绪相对温和，黄=中性/有分歧，红=情绪过热或有明确风险信号，蓝底=HK Anchor特殊行。}
\end{exhibitbox}

读图表7-3须把握三条情绪规律。第一，\textbf{估值越高的标的，公司澄清和监管互动越频繁}：宏昌电子（PE失效）、宏和科技（PE ~750x）、金安国纪（伪AI概念）、中材科技（AI电子布仅2.08\%）均发布过风险提示或澄清公告，且内部人减持与股价高位高度同步——这是典型的"估值泡沫期公司主动降温"信号，表明管理层也认为当前股价偏离基本面。第二，\textbf{融资拥挤度与龙虎榜频次正相关，但与基本面质量负相关}：容大感光融资占比4.87\%（扩充标的中最高）、东材科技4.22\%，均处高位，但两家Q1扣非利润均弱于表观（容大-24\%、东材扣非仅+42\%远低于表观+103\%），说明杠杆资金追逐的是主题热度而非盈利质量。第三，\textbf{30日主力净流几乎全线为负}（仅宏和科技+1.62亿微正），与核心6家呈现的"高位震荡+筹码扩散"格局一致——说明无论是核心还是扩充标的，本轮AI行情都已进入资金面松动阶段，情绪指标领先基本面指标走弱。

Watchlist九家中情绪风险最高的三只依次为：\textbf{金安国纪}（两次澄清无AI敞口但年内涨超500\%，典型伪AI概念泡沫，监管异常波动关注）、\textbf{宏昌电子}（PE失效、7天5涨停、财务负责人减持、公司紧急风险提示四信号共振）、\textbf{宏和科技}（PE ~750x叙事溢价极值、董事长/董秘高位减持、公司主动提示扩产不确定）。这三只均处于"情绪远超基本面"的极值区，一旦主题情绪退潮，估值压缩幅度可能最为剧烈，应作为情绪风险预警的重点监测对象。

\subsection*{Demand Indicator：情绪温度计（2家）}

工业富联（601138）和中际旭创（300308）不直接生产PCB/CCL/铜箔，不作为链内投资标的，但承担"AI算力需求情绪温度计"的角色——二者的股价、融资余额和资金流向是全市场AI算力主题的情绪风向标，其拐点往往领先PCB/CCL链内标的1--2周。

工业富联当前融资余额约105.07亿元、占流通市值仅0.68\%（杠杆拥挤度低），基金持有约995行、合计持股市值约697.91亿元，近90日未登上龙虎榜（大盘蓝筹属性），30日主力净流-43.55亿元但最新单日净流入+45.68亿元呈现剧烈反转。中际旭创融资余额约430.27亿元、占流通市值3.04\%（绝对规模大但占比不极端），基金持有约994行、合计持股市值约1,832.77亿元，近90日同样未登龙虎榜，30日主力净流-50.84亿元但最新单日净流入+24.12亿元。

两只需求指标股的共同特征是：\textbf{融资绝对规模大但占流通市值比例不算极端，基金持仓极其分散（均接近995行上限），30日资金流为负但单日反转剧烈}。这说明AI算力主题已从"全面进攻"进入"高位震荡+分歧加大"阶段——不是所有资金都离场，但增量资金入场节奏明显放缓，且对单个催化（如GB300出货、1.6T光模块季度数据）的反应越来越短促。对PCB/CCL链内标的而言，需求指标股的情绪温度是重要领先信号：若工业富联和中际旭创双双跌破20日均线且融资余额持续下降，意味着AI算力主题进入情绪退潮期，链内高估值标的的估值压缩风险将显著上升。

\subsection*{HK Anchor：港股情绪锚（建滔积层板 01888.HK）}

建滔积层板作为全球CCL龙头和港股通标的，承担两个特殊情绪参照功能：一是全球CCL量价风向标——其涨价/降价决策直接影响A股生益、华正、南亚等CCL厂的盈利预期；二是港股AI-PCB主题的外资情绪窗口——南向资金流向和沽空比例反映了跨市场资金对高端CCL景气的看法。

情绪指标方面，建滔积层板2025H2至2026年累计提价超50\%，股价年内一度飙升约570\%（半年市值暴涨约2,200亿港元），是本轮全球AI-CCL景气的最强价格信号。港股通南向资金连续8日净买入累计96.53亿港元，显示A/H两地资金对CCL涨价逻辑的高度共识。卖方目标价区间（国信82.33--89.39港元、花旗120港元）与当前股价形成上行动能参照，但花旗120港元较国信区间上沿仍有约34\%的差距，说明外资内部对涨价周期延续时长也存在分歧。

\textbf{港股锚的预警价值大于进攻价值}。若建滔积层板股价见顶回落、或公司宣布涨价暂停/降价，将是A股CCL/PCB板块估值拥挤度的最直接预警信号——因为建滔作为纯CCL标的，其定价最直接地反映了CCL周期的边际变化，不受A股特有的主题炒作和流动性溢价干扰。对本报告持仓而言，建滔的量价信号是材料链仓位调整的外生验证锚，而非直接交易标的。

扩充标的情绪画像的整体结论是：\textbf{14只扩充标的的情绪热度分层显著，与基本面质量呈反向关系}——Core Candidate三家中景旺情绪相对克制、东山和芯碁偏热；Watchlist九家情绪风险集中在宏昌/宏和/金安三只"伪AI+高估值"标的；Demand Indicator显示AI算力主题已从全面进攻进入高位震荡分歧期；HK Anchor确认CCL涨价周期仍在延续但外资内部分歧加大。整体情绪温度读数为"偏热、边际松动"，与核心6家呈现的"高位震荡+筹码扩散"格局相互印证。
